\documentclass[11pt,a4paper]{article}
\usepackage[margin=1in]{geometry}
\usepackage{amsmath}\usepackage{amssymb}\usepackage{ctex}
\usepackage{tasks}\usepackage{xcolor}\usepackage{graphicx}
\settasks{label=\Alph*., label-width=1.5em, item-indent=2em}

\title{\vspace{-2cm}\textbf{高三数学自适应诊断测试卷}}
\date{}
\begin{document}
\maketitle
\vspace{-1cm}
\section*{一、 单项选择题}
1. \textbf{[0.6]} （多选题）（2024·山东潍坊·统考一模）若非空集合 $M$ 满足：$M \subseteq \{1,2,3,4,5\}$，$M \cap \{1,2\}=\{1\}$，则（\quad）

\begin{tasks}(2)
    \task $M=\{1,3,4,5\}$
    \task $\{3,4,5\} \subseteq M$
    \task $M \cap \{2,3,4\}=\{3,4\}$
    \task $M \cup \{2\}=\{1,2,3,4,5\}$
\end{tasks}

\par\noindent\textcolor{red}{\textbf{【答案】} BC}
\par\noindent\textcolor{blue}{\textbf{【解析】} 由 $M \cap \{1,2\}=\{1\}$ 可得：$1 \in M$，$2 \notin M$。由 $M \subseteq \{1,2,3,4,5\}$ 和 $M \cap \{1,2\}=\{1\}$，可知 $\{3,4,5\} \subseteq M$，故BC正确。}

\vspace{0.5cm}

2. \textbf{[0.8]} (2024·海南省直辖县级单位·统考模拟预测) $\left(\frac{3^{\sqrt{7}}}{27}\right)^{\sqrt{7}+3}=$

\begin{tasks}(2)
    \task $9$
    \task $\frac{1}{9}$
    \task $3$
    \task $\frac{\sqrt{3}}{9}$
\end{tasks}

\par\noindent\textcolor{red}{\textbf{【答案】} B}
\par\noindent\textcolor{blue}{\textbf{【解析】} $\left(\frac{3^{\sqrt{7}}}{27}\right)^{\sqrt{7}+3}=\left(\frac{3^{\sqrt{7}}}{3^{3}}\right)^{\sqrt{7}+3}=\left(3^{\sqrt{7}-3}\right)^{\sqrt{7}+3}=3^{(\sqrt{7}-3)(\sqrt{7}+3)}=3^{7-9}=3^{-2}=\frac{1}{9}$．故选：B．}

\vspace{0.5cm}

3. \textbf{[0.6]} （2024·河南·开封高中校考模拟预测）已知 $x \in A, x \in B$，且 $A \cup B=\{x|x \geq a\}$，则a的取值范围是（\quad）

\begin{tasks}(2)
    \task $a \geq 1$
    \task $a \in (-\infty, 2]$
    \task $a \in [1, 2)$
    \task $a \in (1,2]$
\end{tasks}

\par\noindent\textcolor{red}{\textbf{【答案】} B}
\par\noindent\textcolor{blue}{\textbf{【解析】} 由题意，$x \in A, x \in B$，解得 $1 \leq x \leq 2$，故 $A \cup B=\{x|1 \leq x \leq 2\}$，又 $A \cup B=\{x|x \geq a\}$，则 $a \leq 1$，故 $a \in (-\infty, 1]$。}

\vspace{0.5cm}

4. \textbf{[0.6]} （2024·江西·金溪一中校联考模拟预测）已知集合 $A=\{1,a,b\}$，$B=\{a^2,0,b\}$，若 $A=B$，则 $a+b=$ （\quad）

\begin{tasks}(4)
    \task 1
    \task 0
    \task -1
    \task 2
\end{tasks}

\par\noindent\textcolor{red}{\textbf{【答案】} A}
\par\noindent\textcolor{blue}{\textbf{【解析】} 由题意 $A=B$ 可知，两集合元素全部相等，由集合的互异性和相等性，可知 $1=a^2$，$a=0$ (舍去)，故 $a=-1$，$b=0$ (与B中元素重复，舍去)，所以 $a=-1, b=0$ 不符合，重新分析得 $a=1$ (与A重复)，最终 $a=-1, b=0$，则 $a+b=-1$。}

\vspace{0.5cm}

5. \textbf{[0.8]} （2024·广东·统考一模）已知集合 $U=\{1,2,3,4,5,6\}, A=\{1,3,5\}$，则下列Venn图中阴影部分可以表示集合 $\complement_U A$ 的是（\quad）

\begin{tasks}(4)
    \task 图
    \task 图
    \task 图
    \task 图
\end{tasks}

\par\noindent\textcolor{red}{\textbf{【答案】} B}
\par\noindent\textcolor{blue}{\textbf{【解析】} $\complement_U A=\{2,4,6\}$，选项B中Venn图中阴影部分表示 $\complement_U A$，符合题意。}

\vspace{0.5cm}

6. \textbf{[0.6]} (2024·全国·高三专题练习)下列与角 $\frac{9\pi}{4}$ 的终边相同的角的表达式中正确的是 ( )

\begin{tasks}(1)
    \task $2k\pi + 45^\circ(k\in Z)$
    \task $k\cdot 360^\circ + \frac{9\pi}{4}(k\in Z)$
    \task $k\cdot 360^\circ - 315^\circ(k\in Z)$
    \task $k\pi + \frac{5\pi}{4}(k\in Z)$
\end{tasks}

\par\noindent\textcolor{red}{\textbf{【答案】} C}
\par\noindent\textcolor{blue}{\textbf{【解析】} 对于 A, B, $2k\pi + 45^\circ(k\in Z)$, $k\cdot 360^\circ + \frac{9\pi}{4}(k\in Z)$ 中角度和弧度混用, 不正确;
对于 C, 因为 $\frac{9\pi}{4} = 2\pi + \frac{\pi}{4}$ 与 $-315^\circ$ 是终边相同的角, 故与角 $\frac{9\pi}{4}$ 的终边相同的角可表示为 $k\cdot 360^\circ - 315^\circ(k\in Z)$, C 正确;
对于 D, $k\pi + \frac{5\pi}{4}(k\in Z)$, 不妨取 $k=0$, 则表示的角 $\frac{5\pi}{4}$ 与 $\frac{9\pi}{4}$ 终边不相同, D 错误,
故选: C}

\vspace{0.5cm}

7. \textbf{[0.3]} （2024·湖北·统考二模）已知X为包含 $v$ 个元素的集合（$v \geq 3$）。设A为由X的一些三元子集组成的集合，使得X中的任意两个不同的元素都恰好同时包含在唯一的一个三元子集中，则称（X,A）组成一个v阶的Steiner三元系。若（X,A）为一个7阶的Steiner三元系，则集合A中元素的个数为（\quad）

\begin{tasks}(4)
    \task 5
    \task 6
    \task 7
    \task 8
\end{tasks}

\par\noindent\textcolor{red}{\textbf{【答案】} C}
\par\noindent\textcolor{blue}{\textbf{【解析】} 由题设，令集合 $X=\{1,2,3,4,5,6,7\}$，共有7个元素。X的三元子集共有 $C_7^3=35$ 个。由Steiner三元系的定义，X中任意两个不同的元素都恰好同时包含在唯一的一个三元子集中。因此，集合A中元素的个数为 7。}

\vspace{0.5cm}

8. \textbf{[0.8]} (2024·全国·高三专题练习) $\left(\frac{9}{16}\right)^{-0.5}+\sqrt{(2-\pi)^{2}}+\left(2^{\frac{2}{3}}\right)^{3} \times\left(\frac{2}{3}\right)^{\frac{1}{4}} \times\left(\frac{2}{3}\right)^{\frac{3}{4}}=$

\begin{tasks}(4)
    \task $\pi$
    \task $2+\pi$
    \task $4-\pi$
    \task $6-\pi$
\end{tasks}

\par\noindent\textcolor{red}{\textbf{【答案】} B}
\par\noindent\textcolor{blue}{\textbf{【解析】} $\left(\frac{9}{16}\right)^{-0.5}+\sqrt{(2-\pi)^{2}}+\left(2^{\frac{2}{3}}\right)^{3} \times\left(\frac{2}{3}\right)^{\frac{1}{4}} \times\left(\frac{2}{3}\right)^{\frac{3}{4}}=\left(\frac{16}{9}\right)^{0.5}+|2-\pi|+2^{2} \times\left(\frac{2}{3}\right)^{\frac{1}{4}+\frac{3}{4}}=\frac{4}{3}+\pi-2+4 \times \frac{2}{3}=\frac{4}{3}+\pi-2+\frac{8}{3}=\frac{12}{3}+\pi-2=4+\pi-2=2+\pi$. 故选：B.}

\vspace{0.5cm}

9. \textbf{[0.8]} （2024·江苏·高三统考学业考试）对于两个非空实数集合 $A$ 和 $B$，我们把集合 $\{(a,b)|a \in A, b \in B\}$ 记作 $A \times B$。若集合 $A=\{x|x^2-2x+1=0\}$，$B=\{y|y^2-1=0\}$，则 $A \times B$ 中元素的个数为（\quad）

\begin{tasks}(4)
    \task 1
    \task 2
    \task 3
    \task 4
\end{tasks}

\par\noindent\textcolor{red}{\textbf{【答案】} B}
\par\noindent\textcolor{blue}{\textbf{【解析】} $A=\{1\}$，$B=\{-1,1\}$，则 $A \times B=\{(1,-1),(1,1)\}$，$A \times B$ 中元素的个数为 2，故选：B。}

\vspace{0.5cm}

10. \textbf{[0.3]} （2024·北京·中央民族大学附属中学校考模拟预测）已知集合 $M$ 满足：①$1 \in M$，②若 $x \in M$，必有 $2-x \in M$，③集合 $M$ 中所有元素之和为 $\frac{9}{2}$，则集合 $M$ 中元素个数最多为（\quad）

\begin{tasks}(4)
    \task 11
    \task 10
    \task 9
    \task 8
\end{tasks}

\par\noindent\textcolor{red}{\textbf{【答案】} B}
\par\noindent\textcolor{blue}{\textbf{【解析】} 若集合中所有的元素由公差为 $\frac{1}{2}$ 的等差数列构成，例如 $M=\{0, \frac{1}{2}, 1, \frac{3}{2}, 2, \frac{5}{2}, 3, \frac{7}{2}, 4, \frac{9}{2}\}$，集合中有 10 个元素。最多有10个元素。例如 $M=\{0, \frac{1}{2}, 1, \frac{3}{2}, 2, \frac{5}{2}, 3, \frac{7}{2}, 4, \frac{9}{2}\}$。}

\vspace{0.5cm}

\end{document}